\chapter{The no-go constitution}\label{chap:nogo}

The theory must distinguish a genuine comparison map from a coincidence of
shadows.  Each prohibition below names the datum whose absence blocks the
proposed identification and the construction that would make a comparison
possible.

\section{Carrier prohibitions}
\begin{enumerate}[series=nogo,label=\textbf{N\arabic*.},ref=N\arabic*]
\item \textbf{Residue is not the transverse bar differential.}  Residue changes normal/collision degree; the bar differential contracts adjacent transverse letters.  They can anticommute in a bicomplex.
\item \textbf{A logarithmic form does not extract all principal parts.}  Higher poles require the complete normal-jet filtration or a chosen connection.
\item \textbf{Incidence does not contract screen cohomology.}  A terminal corolla can make an incidence nerve contractible while the geometric screen has nontrivial cohomology.
\item \textbf{Order does not produce exchange.}  An ordered chamber on a line is contractible and has no braid generator.
\item \textbf{A trace is not an augmentation.}  Cyclic linear functionals do not define multiplicative bar faces.
\item \textbf{Graph topology does not determine coefficients.}  A stable graph specifies possible contractions; a pairing and state assignment specify their values.
\end{enumerate}

\section{Monoidal and operadic prohibitions}
\begin{enumerate}[resume*=nogo,label=\textbf{N\arabic*.}]
\item \textbf{Hadamard and Cauchy products are not globally strongly duoidal.}  The exact relation is a split inclusion/projection with aligned idempotent.
\item \textbf{Transverse \(\Eone\) does not imply associative chiral \(\Eone\).}  The two products have different carriers and operads.
\item \textbf{Associative order does not imply a quantum vertex \(\mathcal S\)-operator.}  Unitarity, Yang--Baxter, shift, and braided locality are extra equations.
\item \textbf{Two pairwise exchange rules do not imply coherent three-field exchange.}  The critical-pair obstruction is Yang--Baxter.
\item \textbf{The interval bar, associative chiral bar, and chiral Chevalley complex are not one complex.}  A comparison requires an enveloping or distributive-law map.
\item \textbf{A quadratic dual is not automatically the derived Koszul dual.}  Exactness of the quadratic Koszul complex is the missing theorem.
\item \textbf{Bar--cobar is not an equivalence without completion.}  The output is nilcompletion/conilcompletion in general.
\item \textbf{Verdier, boundary, quadratic, and bar duals are not definitionally identical.}  Finite self-dual examples supply comparison maps.
\end{enumerate}

\section{Algebra and calculation prohibitions}
\begin{enumerate}[resume*=nogo,label=\textbf{N\arabic*.}]
\item \textbf{The Weyl algebra has no scalar augmentation.}  Use a Rees degeneration, a curved dual, or a relative module.
\item \textbf{A full matrix algebra has no scalar augmentation.}  Use Morita/relative bar data or embed it in an honestly augmented algebra.
\item \textbf{A multi-idempotent quiver is naturally relative over its vertex algebra.}  Scalar augmentation forgets the boundary objects.
\item \textbf{A chain-space count is not a homology dimension.}  The differential and its rank are part of the invariant.
\item \textbf{A matching integer sequence is not a quasi-isomorphism.}  Motzkin and Riordan arithmetic survives only in its actual planar combinatorial role.
\item \textbf{Central charge does not classify a chiral algebra.}  OPEs, modules, fusion, modular data, and transverse multiplication remain independent.
\item \textbf{A character does not reconstruct a Hopf algebra.}  Coproduct, antipode, associator, and exchange are structured maps.
\item \textbf{A denominator product does not construct a Hall realization.}  Moduli, orientations, extension correspondences, and PBW comparison are separate.
\end{enumerate}

\section{Physics and sewing prohibitions}
\begin{enumerate}[resume*=nogo,label=\textbf{N\arabic*.}]
\item \textbf{A boundary centre does not determine an arbitrary decoupled bulk.}  An open--closed equivalence is a theorem with generation and locality hypotheses.
\item \textbf{The two genera are not universally equal.}  Ribbon and Deligne--Mumford handles are different edge colours.
\item \textbf{Formal sewing does not imply analytic convergence.}  Energy bounds or trace-class estimates are additional inputs.
\item \textbf{Different anomaly classes are not one scalar.}  Curved, BV, BRST, determinant, and framing anomalies live in different complexes.
\item \textbf{Classical factorisation does not guarantee quantum factorisation.}  The QME obstruction and boundary anomaly must vanish.
\item \textbf{A physical analogy does not supply a functor.}  Every comparison between gauge theory, Hall geometry, vertex algebra, and automorphic form must specify source, target, and map.
\end{enumerate}

\section{Positive replacements}
The constitution is not merely restrictive.  Each no-go theorem has a constructive replacement:
\begin{center}
\begin{tabularx}{0.96\textwidth}{p{0.31\textwidth}X}
\toprule
False shortcut & Correct construction\\
\midrule
residue as bar differential & total complex with typed collision and bar degrees\\
strong duoidal equivalence & aligned-decomposition retract and locality squares\\
``nonlocal'' from transverse order & independent associative chiral operad and formal-disc recognition\\
ordinary quotient by relations & homotopy pushout of free cofibrant algebras\\
bar of a Weyl/matrix algebra & curved, Rees, relative, or Morita bar\\
sequence recognition & explicit chain matrices and quasi-isomorphism\\
boundary centre equals bulk & bulk-to-centre map plus open--closed theorem\\
one genus & bicoloured stable graphs and mixed recognition\\
one anomaly number & comparison maps between typed deformation complexes\\
automorphic product as algebra & bar character theorem plus separate Hall realization\\
\bottomrule
\end{tabularx}
\end{center}

\chapter{Typed obstruction and finite-calculation ledger}\label{chap:corpus}

\section{Carrier-to-construction atlas}
The mathematical claims separate according to their carriers.
\begin{longtable}{p{0.24\textwidth}p{0.31\textwidth}p{0.34\textwidth}}
\toprule
\textbf{Carrier} & \textbf{Available construction} & \textbf{Additional datum required}\\
\midrule
\endfirsthead
\toprule
\textbf{Carrier} & \textbf{Available construction} & \textbf{Additional datum required}\\
\midrule
\endhead
factorisable \(D\)-modules & aligned pointwise/chiral locality & a mixed distributive law off the aligned sector\\
augmented transverse algebra & interval bar and its word differential & a separate chiral multiplication\\
associative chiral algebra & collision-tree operadic bar & an enveloping or distributive-law comparison with the interval bar\\
relative quiver algebra over \(S\) & \(S\otimes_A^{\mathbb L}S\) & no scalar augmentation is inferred\\
exchange algebra & KZ, ZF, or RTT transport & flatness, unitarity, Yang--Baxter, and a fibre functor as applicable\\
cyclic or modular object & trace and sewing operations & pairing, convergence, and anomaly data\\
root-graded bar complex & Euler supercharacter denominator & a separate Hall correspondence and PBW comparison\\
\bottomrule
\end{longtable}

\section{Finite calculations and their boundaries}
For the quantum plane at \(q=2\) and the Jordan plane, normalized bar
words are graded by total polynomial weight \(w\leq5\) and word length.
The alternating adjacent-multiplication matrices are rational.  Their
nonzero homology in this window is
\[
 H_{1,1}\cong\kfield^2,\qquad H_{2,2}\cong\kfield,
\]
with all other groups zero; Appendix~\ref{app:computations} gives the
bases, matrix ranks, and cycles.  For \(U(\mathfrak{sl}_2)\), the
Chevalley--Eilenberg differential has \(\det d_2=-4\), \(d_3=0\), and
\[
(\dim H_0,\dim H_1,\dim H_2,\dim H_3)=(1,0,0,1).
\]
These calculations refute universal Motzkin or Riordan formulas for the
corresponding ordered bar homology.  A scalar sequence, chain-space count,
or finite window does not establish a quasi-isomorphism outside its stated
carrier and range.

\section{The exact evidentiary standard}
A worked example is complete only when it contains:
\begin{enumerate}
\item the coefficient object and its diagonal extension;
\item the associative product and augmentation or relative base;
\item generators, relations, and a normal-form or PBW theorem;
\item the actual bar or Koszul differential;
\item cycles, boundaries, and homology, not only chain dimensions;
\item any exchange, trace, modular, or physical realization as a named extra map;
\item an exact finite-window calculation, with its parameter range, when constants or signs are nontrivial.
\end{enumerate}
The quantum plane, Jordan plane, relation quiver, \(\mathfrak{sl}_2\), ZF system, \(A_2\) Hall algebra, stable trace bracket, and two-genus Frobenius/CohFT model meet this standard in the present volume.

\chapter{The definitive structure theorem}\label{chap:structure}

\section{The primitive datum}
Let \(X\) be a smooth complex curve.  The complete local object is the tuple
\[
 \mathbb A=
 \bigl(F,m_\perp,m_X,\lambda,\eps;\ \text{completions}\bigr),
\]
where:
\begin{enumerate}
\item \(F\) is a factorisable \(D\)-module carrying every diagonal singularity and regular product;
\item \(m_\perp\) is an augmented transverse \(\Eone\)-multiplication for \(\tensorpt\);
\item \(m_X\) is an augmented associative chiral \(\Eone\)-multiplication for \(\starc\) or \(\Assch\);
\item \(\lambda\) is the aligned mixed distributive law, including all higher homotopies;
\item \(\eps\) is a scalar or relative augmentation appropriate to the algebra;
\item the nilpotent, weight, pole, coradical, and analytic completions are specified separately.
\end{enumerate}
Any subset of this tuple defines the corresponding partial theory.  No missing component is inferred from the others.

\begin{theorem}[Definitive structure theorem]\label{thm:definitive}
For a complete datum \(\mathbb A\) as above, the following constructions are functorial and type preserving.
\begin{enumerate}
\item \(\Barop^\perp(\mathbb A)\) is interval factorisation homology and a conilpotent transverse coalgebra retaining the associative chiral coefficient.
\item \(\Barop^{\ch}_{\Ass}(\mathbb A)\) is the operadic associative chiral bar retaining the transverse algebra.
\item The mixed distributive law gives a total two-direction bar with anticommuting edge differentials and an aligned Fubini comparison.
\item On the nilcomplete/conilcomplete locus, the corresponding bar and cobar functors are inverse.
\item The augmentation boundary, relative module categories, two-sided twisting transform, and derived centre form the Koszul--Morita atlas.
\item An \(\Etwo\)-lift or flat meromorphic transport adds exchange; a fibre functor reconstructs the corresponding quantum symmetry.
\item A compatible BRST operator forms a bicomplex and descends all commuting structures to cohomology.
\item A cyclic transverse pairing and a modular chiral pairing satisfying the mixed square produce the two-edge master equation with independent genus variables.
\item A perturbative BV realization gives the preceding object when renormalised graph weights extend, the QME holds, and the observables algebraize.
\item Root-graded Euler character, trace, period, index, and monodromy maps produce scalar shadows only after the structured object has been constructed.
\end{enumerate}
Conversely, every construction assembled from the displayed carriers factors
through one or more of these functors.  A proposed construction fails whenever
it omits one of the required data or hypotheses.
\end{theorem}
\begin{proof}
Items 1--3 are \cref{thm:bar-interval,thm:two-bars}; item 4 is
\cref{thm:internal-kd,thm:positive-reconstruction}; item 5 is
\cref{thm:double-centralizer,thm:module-comodule}; item 6 follows from
\cref{prop:kz-flat,thm:zf-diamond} and Tannaka reconstruction; item 7 is
\cref{thm:bar-brst}; item 8 is \cref{thm:two-edge-recognition}, realized by
\cref{thm:realized-two-edge}; item 9 is \cref{thm:feynman-weight}; item 10
is exemplified by \cref{thm:stable-trace,thm:bar-denominator}.  Conversely,
each operation factors by type as collision, ordered composition, exchange,
trace, sewing, reduction, pushforward, or scalar evaluation, and each appears
in the atlas.  The no-go constitution records the missing input in every
forbidden identification.
\end{proof}

\section{The single invariant sentence}
The complete statement is:
\begin{quote}
A genuinely noncommutative chiral field theory on a curve may carry two independent associative structures: ordered fusion in a transverse topological direction and associative, possibly nonlocal, chiral composition on the curve.  Their coherent coexistence is an algebra over a two-direction cofibrant operad, with locality supported on the aligned-decomposition sector.  Bar, centre, exchange, trace, sewing, reduction, and arithmetic are functorial constructions from this typed datum, never substitutes for it.
\end{quote}

\section{Consequences of the typed datum}
The typed datum has four consequences.
\begin{enumerate}
\item The aligned projector is the organizing law of mixed locality.
\item It identifies the missing associative chiral/nonlocal \(\Eone\) as a second operadic direction rather than a synonym for transverse order.
\item It turns generators and relations into a derived, calculable presentation theory over the two-coloured cofibrant operad.
\item The examples are chain-level calculations that can falsify a proposed identification as well as illustrate it.
\end{enumerate}

\section{Coda: the geometry remembers}
The diagonal remembers singularity.  The interval remembers order.  The formal disc remembers associative chiral iteration.  The plane remembers exchange.  The circle remembers trace.  The stable curve remembers sewing.  The compactified boundary remembers the equations.  A scalar remembers only the value of a chosen functional.  The typed theory keeps these memories separate until a theorem joins them.
